\documentclass[12pt,a4paper]{article}
\usepackage[utf8]{inputenc}
\usepackage[T1]{fontenc}
\usepackage{lmodern}
\usepackage{hyperref}
\usepackage{graphicx}
\usepackage{amsmath}
\usepackage{amssymb}
\usepackage{listings}
\usepackage{xcolor}
\usepackage{booktabs}
\usepackage{geometry}
\usepackage{fancyhdr}
\usepackage{titlesec}
\usepackage{enumitem}

% Page geometry
\geometry{a4paper, margin=0.8in}

% Header and footer
\pagestyle{fancy}
\fancyhf{}
\fancyhead[L]{Conference Management System}
\fancyhead[R]{DBMS Project}
\fancyfoot[C]{\thepage}
\renewcommand{\headrulewidth}{0.4pt}
\renewcommand{\footrulewidth}{0.4pt}

% Set the headheight to avoid warnings
\setlength{\headheight}{15pt}

% Title formatting
\titleformat{\section}{\normalfont\Large\bfseries}{\thesection}{1em}{}
\titleformat{\subsection}{\normalfont\large\bfseries}{\thesubsection}{1em}{}

% Code listing style
\definecolor{codegreen}{rgb}{0,0.6,0}
\definecolor{codegray}{rgb}{0.5,0.5,0.5}
\definecolor{codepurple}{rgb}{0.58,0,0.82}
\definecolor{backcolour}{rgb}{0.95,0.95,0.92}

\lstdefinestyle{mystyle}{
    backgroundcolor=\color{backcolour},   
    commentstyle=\color{codegreen},
    keywordstyle=\color{codepurple},
    stringstyle=\color{codegreen},
    basicstyle=\ttfamily\footnotesize,
    breakatwhitespace=false,         
    breaklines=true,                 
    captionpos=b,                    
    keepspaces=true,                 
    numbers=left,                    
    numbersep=5pt,                  
    showspaces=false,                
    showstringspaces=false,
    showtabs=false,                  
    tabsize=2
}

% Define JavaScript language for listings
\lstdefinelanguage{JavaScript}{
  keywords={async, await, break, case, catch, class, const, continue, debugger, default, delete, do, else, export, extends, false, finally, for, function, if, import, in, instanceof, new, null, return, super, switch, this, throw, true, try, typeof, var, void, while, with, yield, let},
  keywordstyle=\color{codepurple}\bfseries,
  ndkeywords={boolean, number, string, object, array},
  ndkeywordstyle=\color{codegray}\bfseries,
  identifierstyle=\color{black},
  sensitive=false,
  comment=[l]{//},
  morecomment=[s]{/*}{*/},
  commentstyle=\color{codegreen},
  stringstyle=\color{codegreen},
  morestring=[b]',
  morestring=[b]"
}

\lstset{style=mystyle}

\begin{document}

\begin{titlepage}
    \centering
    \vspace*{1cm}
    {\scshape\LARGE Indian Institute of Technology Bombay \par}
    \vspace{1.5cm}
    {\huge\bfseries Conference Management System \par}
    \vspace{2cm}
    {\Large\itshape DBMS Project Report \par}
    \vspace{1.5cm}
    {\large Prepared by: \par}
    \vspace{0.5cm}
    {Piyush Babar (22B3909)\par}
    {Aryan Gupta (22B2255)\par}
    {Ashish Zantye (22B2217)\par}
    {Shresth Verma (22B2211)\par}
    \vfill
    {\large \today \par}
\end{titlepage}

% \tableofcontents
% \newpage

\section{Introduction}
\subsection{Project Overview}
The Conference Management System is a comprehensive web application designed to assist academic conference organizers in managing the peer review process. The system focuses on detecting potential conflicts of interest between authors and reviewers, ensuring fair and unbiased paper evaluations.

\subsection{Objectives}
The primary objectives of this project are:
\begin{itemize}
    \item Create a robust database schema for storing academic publication data
    \item Implement efficient data loading from the DBLP computer science bibliography
    \item Develop algorithms for detecting conflicts of interest between authors and reviewers
    \item Create a user-friendly web interface for interacting with the system
    \item Implement features for exploring and visualizing the publication database
    \item Develop a recommendation system for reviewer suggestions accounting for conflicts of interest, reviewer expertise and diversity
\end{itemize}

\subsection{Technologies Used}
The project utilizes the following technologies:
\begin{itemize}
    \item \textbf{Database}: PostgreSQL for data storage and retrieval
    \item \textbf{Backend}: Node.js with Express for server-side logic
    \item \textbf{Frontend}: HTML5, CSS3, and vanilla JavaScript
    \item \textbf{Data Processing}: Python for data loading and preprocessing
    \item \textbf{External APIs}: OpenAlex API for institutional affiliation data
\end{itemize}

\subsection{Project Scope}
The system focuses on two main types of conflict detection:
\begin{enumerate}
    \item \textbf{Coauthorship Conflicts}: When authors and reviewers have collaborated on publications
    \item \textbf{Institutional Conflicts}: When authors and reviewers are affiliated with the same institutions
\end{enumerate}

Additionally, the system provides data exploration features and reviewer recommendation functionality to assist conference organizers.

\newpage
\section{Data Loading}

\subsection{Database Schema Design}
PostgreSQL database schema for DBLP bibliographic data consists of:

\begin{itemize}
    \item \textbf{publications}: Primary table with key fields (title, year, journal, publisher)
    \item \textbf{persons}: Author/editor information storage
    \item \textbf{publication\_authors}: Many-to-many relationship junction table
    \item \textbf{publication\_editors}: Junction table for editors
    \item \textbf{publication\_citations}: Citation relationships
    \item \textbf{author\_affiliations}: Institutional affiliations from OpenAlex
\end{itemize}

Indexes implemented on key lookup fields: author names, publication years, titles, and journals. The complete schema is given at the end of the report.

\subsection{DBLP Data Loading Implementation}
Key implementation components in \texttt{dblp\_load.py}:

\begin{itemize}
    \item XML parsing with lxml's iterparse for memory efficiency (O(1) memory usage)
    \item Transaction management with psycopg2 connection pooling
    \item Custom data type handling for dates, integers and null values
    \item Automatic primary/foreign key management with SQL's ON CONFLICT directives
\end{itemize}

Processing sequence:
1. Connect to PostgreSQL with connection pooling
2. Parse XML stream with iterparse
3. Extract publication metadata
4. Insert publication record
5. Extract and insert authors/editors with relationships
6. Insert citation data
7. Commit transaction batch

\subsection{OpenAlex Affiliation Integration}
Implementation details of \texttt{openalex\_affiliations.py}:

\begin{itemize}
    \item DOI extraction from publication URLs using regex pattern matching
    \item Batch processing with configurable BATCH\_SIZE constant
    \item HTTP request rate limiting (1 req/second) to comply with API terms
    \item JSON response parsing with error handling
    \item Checkpoint mechanism for resume capability
\end{itemize}

API integration flow:
1. Query database for publications with DOIs
2. Extract clean DOIs from URLs
3. Fetch author affiliation data from OpenAlex API
4. Process response and extract institution information
5. Store normalized affiliation data
6. Save checkpoint after each batch
\newpage

\section{Conflict Detection Algorithms}
\subsection{Coauthorship Detection SQL Implementation}
\textbf{Recent coauthorship SQL query:}
\begin{verbatim}
SELECT pa.person_key AS author_id, 
       pr.person_key AS reviewer_id, 
       pa.publication_key
FROM publication_authors pa
JOIN publication_authors pr 
     ON pr.publication_key = pa.publication_key
JOIN publications p 
     ON pa.publication_key = p.key
WHERE pa.person_key = ANY($1::text[])
  AND pr.person_key = ANY($2::text[])
  AND p.year >= EXTRACT(YEAR FROM CURRENT_DATE) - $3
\end{verbatim}

\textbf{Frequent coauthorship SQL query:}
\begin{verbatim}
SELECT pa.person_key AS author_id, 
       pr.person_key AS reviewer_id, 
       COUNT(pa.publication_key) AS count
FROM publication_authors pa
JOIN publication_authors pr 
     ON pr.publication_key = pa.publication_key
JOIN publications p 
     ON pa.publication_key = p.key
WHERE pa.person_key = ANY($1::text[])
  AND pr.person_key = ANY($2::text[])
  AND p.year >= EXTRACT(YEAR FROM CURRENT_DATE) - $3
GROUP BY pa.person_key, pr.person_key
HAVING COUNT(pa.publication_key) >= $4
\end{verbatim}

\subsection{Institutional Conflict Detection}
\textbf{Institutional conflict SQL query:}
\begin{verbatim}
SELECT a_aff.author_key AS author_id, 
       r_aff.author_key AS reviewer_id, 
       a_aff.institution_name
FROM author_affiliations a_aff
JOIN author_affiliations r_aff 
     ON LOWER(a_aff.institution_name) = 
        LOWER(r_aff.institution_name)
WHERE a_aff.author_key = ANY($1::text[])
  AND r_aff.author_key = ANY($2::text[])
  AND a_aff.author_key != r_aff.author_key
\end{verbatim} 

\section{Reviewer Recommendation Algorithm}
\subsection{Recommendation Process Overview}
The reviewer recommendation system uses a multi-step process to identify suitable reviewers while avoiding conflicts of interest:

\begin{enumerate}
    \item Query potential reviewers by publication count (expertise proxy)
    \item Apply conflict filtering (coauthorship, institutional)
    \item Collect metadata for scoring and diversity analysis
    \item Apply a multi-component scoring algorithm
    \item Use a hybrid selection strategy balancing score and diversity
    \item Return top-N candidates with detailed metadata
\end{enumerate}

\subsection{Conflict Filtering}
The system identifies and excludes three types of conflicts:

\begin{itemize}
    \item \textbf{Recent co-authorship:} Reviewers who co-authored papers with submission authors within a configurable time frame
    \item \textbf{Frequent co-authorship:} Reviewers with more than a threshold number of co-authored papers with submission authors
    \item \textbf{Institutional conflicts:} Reviewers affiliated with the same institutions as submission authors
\end{itemize}

\subsection{Scoring Methodology}
For each eligible reviewer, the algorithm calculates a composite score with the following components:

\begin{enumerate}
    \item \textbf{Base Expertise Score (0-5 points):} Based on publication count
    \begin{verbatim}
    // Base score from publication count (experience)
    let score = Math.min(reviewer.publication_count / 10, 5);
    \end{verbatim}
    
    \item \textbf{Venue Diversity Score:} Rewards reviewers who publish in venues that are uncommon among other potential reviewers
    \begin{verbatim}
    // Venues that are less common get higher scores
    const venueRarity = 1 / (venueFrequency[venue.journal] || 1);
    venueDiversityScore += venueRarity * 2;
    \end{verbatim}
    
    \item \textbf{Institutional Diversity Score:} Rewards reviewers from institutions that are less represented
    \begin{verbatim}
    // Less common institutions get higher scores
    const instRarity = 1 / (institutionFrequency[inst.institution_name] || 1);
    institutionDiversityScore += instRarity * 2;
    \end{verbatim}
    
    \item \textbf{Experience Score (0-3 points):} Based on years active in publishing, with square root applied for diminishing returns
    \begin{verbatim}
    const yearsActive = reviewer.years_active.last_year - 
                        reviewer.years_active.first_year + 1;
    score += Math.min(Math.sqrt(yearsActive) / 2, 3);
    \end{verbatim}
    
    \item \textbf{Collaboration Network Score (0-2 points):} Based on number of co-authors, with logarithmic scaling
    \begin{verbatim}
    // Broader networks indicate greater field awareness
    score += Math.min(Math.log10(reviewer.coauthor_count || 1), 2);
    \end{verbatim}
\end{enumerate}

\subsection{Key Metadata Queries}
The system utilizes several specialized queries to gather metadata for scoring:

\begin{verbatim}
/* Publication venues query */
SELECT DISTINCT p.journal, COUNT(*) as count
FROM publication_authors pa
JOIN publications p ON pa.publication_key = p.key
WHERE pa.person_key = $1 AND p.journal IS NOT NULL
GROUP BY p.journal
ORDER BY count DESC
LIMIT 3

/* Years of activity query */
SELECT MIN(p.year) as first_year, MAX(p.year) as last_year
FROM publication_authors pa
JOIN publications p ON pa.publication_key = p.key
WHERE pa.person_key = $1 AND p.year IS NOT NULL

/* Co-authors count query */
SELECT COUNT(DISTINCT pa2.person_key) as coauthor_count
FROM publication_authors pa1
JOIN publication_authors pa2 ON pa1.publication_key = pa2.publication_key 
                           AND pa1.person_key != pa2.person_key
WHERE pa1.person_key = $1

/* Institutional affiliations query */
SELECT institution_name, city, country, COUNT(*) as frequency
FROM author_affiliations
WHERE author_key = $1
GROUP BY institution_name, city, country
ORDER BY frequency DESC
LIMIT 3
\end{verbatim}

\subsection{Diversity-Aware Selection Strategy}
Rather than simply selecting the top N scored reviewers, the algorithm uses a hybrid strategy:

\begin{enumerate}
    \item First, select the highest-scoring reviewer
    \item Then, alternate between selecting:
    \begin{itemize}
        \item Reviewers who add diversity to the pool (on odd selections)
        \item The next highest-scoring reviewer (on even selections)
    \end{itemize}
\end{enumerate}

The diversity check evaluates whether a reviewer publishes in venues or belongs to institutions not already covered by selected reviewers:

\begin{verbatim}
const addsDiversity = (reviewer) => {
    // If we have fewer than 2 reviewers, any will do
    if (selectedReviewers.length < 2) return true;
    
    // Check for new venues
    for (const venue of reviewer.venues) {
        if (!selectedVenues.has(venue.journal)) {
            return true;
        }
    }
    
    // Check for new institutions
    if (reviewer.institutions) {
        for (const inst of reviewer.institutions) {
            if (!selectedInstitutions.has(inst.institution_name)) {
                return true;
            }
        }
    }
    
    return false;
};
\end{verbatim}

This approach ensures a balance between selecting the most qualified reviewers and maintaining diversity in research focus and institutional backgrounds. 

\newpage
\section{Frontend Implementation}

\subsection{Component Architecture}
Structure of frontend components:
\begin{verbatim}
frontend/
|-- index.html           # Main conflict checking page
|-- view-tables.html     # Database table explorer
|-- recommend-reviewers.html  # Reviewer suggestions
|-- author-profile.html  # Author information display
|-- styles/
    |-- style.css        # Centralized styling
\end{verbatim}

\subsection{Conflict Checking Component}
The conflict checking component provides an intuitive interface for identifying potential conflicts of interest between authors and reviewers. It features dual input fields for author and reviewer IDs with validation that enforces proper formatting through regex pattern matching. The component implements asynchronous form submission to prevent page reloads, providing immediate feedback through a loading indicator during API processing. Upon receiving results, it dynamically generates appropriate data visualizations based on the type and volume of conflicts detected. The system clearly differentiates between coauthorship and institutional conflicts through color-coded indicators and informative tooltips. Error states are handled gracefully with clear user guidance for resolution.

\subsection{Table Viewer Component}
The table viewer component enables database exploration through a flexible interface that adapts to any table structure without requiring predefined templates. It generates table representations dynamically based on metadata received from the backend, automatically adjusting column widths and formatting based on data types. The component implements pagination for large datasets with configurable page sizes to optimize performance. Special data types receive enhanced formatting - timestamps are displayed in human-readable format, URLs become clickable links, and long text fields are truncated with expandable views. The component maintains responsive behavior through CSS class manipulation that adjusts to container dimensions while preserving data readability.

\subsection{Reviewer Recommendation Component}
The reviewer recommendation interface provides a sophisticated system for suggesting appropriate reviewers while avoiding conflicts of interest. It features multiple configuration parameters including conflict lookback periods, collaboration thresholds, and diversity preferences, all with reasonable defaults and contextual guidance. The component presents recommendations in detailed cards showing each reviewer's publication history, venue diversity, institutional affiliations, and experience metrics. A visual score indicator helps users quickly assess reviewer suitability, while expandable sections provide additional detailed information on demand. The interface balances information density with usability through careful visual hierarchy and progressive disclosure of complex data.

\subsection{Author Profile Component}
The author profile component provides comprehensive information about researchers and their publication history. It includes a statistical overview section highlighting publication counts, venue distribution, and collaboration patterns. The publication timeline visualization uses interactive elements to show productivity trends and research focus evolution over time. Institutional affiliation data is presented with geographical context where available. The component also provides collaboration network insights, identifying frequent co-authors and research clusters. Navigation between author profiles is seamless through hyperlinked references, enabling exploration of the broader academic network while maintaining context.

\section{Future Enhancements}
Planned future enhancements for the system include:
\begin{itemize}
    \item Integration of natural language processing for query interpretation
    \item Enhanced visualization of author-reviewer networks
    \item Machine learning for automated reviewer suggestions
    \item Extended conflict detection based on funding sources and personal relationships
    \item Mobile application for on-the-go conference management
\end{itemize}

\section{Conclusion}
The Conference Management System successfully achieves its goal of providing a comprehensive solution for detecting conflicts of interest in the academic peer review process. By combining publication data from DBLP with institutional affiliation information from OpenAlex, the system offers robust conflict detection capabilities.

The project demonstrates the effective use of modern web technologies and database systems to solve a real-world problem in academic publishing. The modular architecture and clean separation of concerns allow for future extensions and enhancements.

The system's user-friendly interface makes it accessible to conference organizers without requiring specialized technical knowledge, while its sophisticated backend algorithms ensure accurate conflict detection and reviewer recommendations.

\newpage
\appendix
\section{Database Structure}

This appendix provides the complete database structure including tables, fields, relationships, and indexes used in the Conference Management System.

\subsection{Core Tables Schema}

\begin{lstlisting}[language=SQL]
-- Publications table - stores all publication types
CREATE TABLE publications (
    key VARCHAR(255) PRIMARY KEY,             -- Unique identifier for the publication
    type VARCHAR(50) NOT NULL,                -- Type of publication ('article', 'book', etc.)
    title TEXT,                               -- Title of the publication
    booktitle TEXT,                           -- Title of the book/proceedings 
    pages VARCHAR(100),                       -- Page range of the publication
    year INTEGER,                             -- Year of publication
    volume VARCHAR(50),                       -- Volume number for journals
    number VARCHAR(50),                       -- Issue number for journals
    month VARCHAR(20),                        -- Month of publication
    journal VARCHAR(255),                     -- Journal name for articles
    publisher VARCHAR(255),                   -- Publisher name
    note TEXT,                                -- Additional notes
    crossref VARCHAR(255),                    -- Cross-reference to another publication
    isbn VARCHAR(50),                         -- ISBN number
    series VARCHAR(255),                      -- Series name for books or proceedings
    school VARCHAR(255),                      -- School name for theses
    chapter VARCHAR(50),                      -- Chapter number for book sections
    publnr VARCHAR(100),                      -- Publisher's identifier
    url TEXT,                                 -- URL pointing to the publication
    ee TEXT,                                  -- Electronic edition URL
    address TEXT,                             -- Publication address/location
    stream TEXT,                              -- Stream information
    rel TEXT,                                 -- Relationship information
    mdate DATE,                               -- Last modification date
    cdate DATE,                               -- Creation date
    publtype VARCHAR(255),                    -- Publication type qualifier
    reviewid VARCHAR(255),                    -- Review identifier
    rating VARCHAR(50),                       -- Publication rating
    cdrom VARCHAR(255)                        -- CD-ROM information
);

-- Persons table - for authors and editors
CREATE TABLE persons (
    key VARCHAR(255) PRIMARY KEY,             -- Unique identifier for the person
    name VARCHAR(255) NOT NULL,               -- Person's name
    mdate DATE,                               -- Last modification date
    cdate DATE                                -- Creation date
);

-- Authors junction table - links publications with authors
CREATE TABLE publication_authors (
    id SERIAL PRIMARY KEY,                    -- Auto-incrementing primary key
    publication_key VARCHAR(255) REFERENCES publications(key) ON DELETE CASCADE,
    person_key VARCHAR(255) REFERENCES persons(key) ON DELETE CASCADE,
    position INTEGER NOT NULL,                -- To maintain order of authors
    aux TEXT,                                 -- Auxiliary information
    bibtex TEXT,                              -- BibTeX information
    orcid VARCHAR(50),                        -- ORCID identifier
    label TEXT,                               -- Label information
    type VARCHAR(100),                        -- Author type
    UNIQUE (publication_key, person_key, position)
);

-- Editors junction table - links publications with editors
CREATE TABLE publication_editors (
    id SERIAL PRIMARY KEY,                    -- Auto-incrementing primary key
    publication_key VARCHAR(255) REFERENCES publications(key) ON DELETE CASCADE,
    person_key VARCHAR(255) REFERENCES persons(key) ON DELETE CASCADE,
    position INTEGER NOT NULL,                -- To maintain order of editors
    aux TEXT,                                 -- Auxiliary information
    orcid VARCHAR(50),                        -- ORCID identifier
    label TEXT,                               -- Label information
    type VARCHAR(100),                        -- Editor type
    UNIQUE (publication_key, person_key, position)
);

-- Citations table - for citations between publications
CREATE TABLE publication_citations (
    id SERIAL PRIMARY KEY,                    -- Auto-incrementing primary key
    publication_key VARCHAR(255) REFERENCES publications(key) ON DELETE CASCADE,
    cite TEXT NOT NULL,                       -- Citation text
    ref VARCHAR(255),                         -- Reference to cited publication
    aux TEXT,                                 -- Auxiliary information
    label TEXT,                               -- Label information
    type VARCHAR(100)                         -- Citation type
);

-- Author affiliations table - for institutional affiliations
CREATE TABLE author_affiliations (
    id SERIAL PRIMARY KEY,                    -- Auto-incrementing primary key
    publication_key TEXT REFERENCES publications(key),
    author_key TEXT REFERENCES persons(key),
    doi TEXT,                                 -- DOI of the publication
    author_name TEXT,                         -- Author name (redundant for performance)
    institution_name TEXT,                    -- Institution/organization name
    city TEXT,                                -- City location
    country TEXT,                             -- Country location
    fetched_at TIMESTAMP DEFAULT CURRENT_TIMESTAMP,
    UNIQUE(publication_key, author_key, institution_name)
);
\end{lstlisting}

\subsection{Indexes}

\begin{lstlisting}[language=SQL]
-- Indexes on publications table
CREATE INDEX idx_publication_type ON publications(type);
CREATE INDEX idx_publication_year ON publications(year);
CREATE INDEX idx_publication_title ON publications(title);
CREATE INDEX idx_publication_journal ON publications(journal);

-- Indexes on persons table
CREATE INDEX idx_person_name ON persons(name);

-- Indexes on junction tables
CREATE INDEX idx_publication_authors_pub ON publication_authors(publication_key);
CREATE INDEX idx_publication_authors_person ON publication_authors(person_key);
CREATE INDEX idx_publication_editors_pub ON publication_editors(publication_key);
CREATE INDEX idx_publication_editors_person ON publication_editors(person_key);

-- Indexes on affiliations table
CREATE INDEX idx_author_affiliations_author ON author_affiliations(author_key);
CREATE INDEX idx_author_affiliations_pub ON author_affiliations(publication_key);
CREATE INDEX idx_author_affiliations_inst ON author_affiliations(institution_name);
\end{lstlisting}

\subsection{Performance Considerations}

The database is optimized for the following common operations:

\begin{itemize}
    \item Finding publications by author (indexed join through publication\_authors)
    \item Finding co-authors of an author (self-join on publication\_authors)
    \item Finding publications by year and type (compound index on publications)
    \item Finding authors by institution (index on author\_affiliations)
    \item Finding institutional conflicts (join on author\_affiliations)
\end{itemize}

\end{document} 